\chapter{Aplicaciones de la termodinámica}

\section{Diagramas de fase de una sustancia pura de un solo componente}

\subsection{Ecuación de Van der Waals}

A pesar de que muchos gases reales obedecen en cierta medida la ecuación de los gases ideales,
muchos de ellos se desvían de dicho comportamiento; para tratar de solventar estas desviaciones, en
$1873$ Johannes Diderik van der Waals postuló una mejora a la ecuación de estado mecánica
($P/T$) [no a la energética] al considerar dos aspectos importantes:
\begin{enumerate}[label=\arabic*.,itemsep=2pt]
  \item Las partículas que forman un gas real \emph{sí} ocupan un espacio, por lo que van der Waals
    propuso una corrección al gas ideal, la cual resta al volumen total el volumen ocupado por las
    partículas que forman este gas.
  \item En un gas real las partículas interactúan entre sí, a través de fuerzas de atracción o
    repulsión; van der Waals introduce una corrección para la presión, de modo que se tomen en
    cuenta dichas fuerzas al medir la presión.
\end{enumerate}
La \textbf{ecuación de estado mecánica de Van der Waals} es entonces
\begin{equation*}
  \boxed{P=\frac{RT}{v-b}-\frac{a}{v^{2}}}
\end{equation*}
donde $a$ y $b$ son las constantes empíricas que se introducen para tales correcciones y que son
características de cada tipo de gas. La constante $b$ corrige el volumen y la constante $a$ corrige
las fuerzas; $V-Nb$ es el volumen corregido, y el término $a/v^{2}$ indica que la interacción entre
partículas disminuye la presión de manera proporcional al número de pares de partículas.

El éxito de esta ecuación es, cuantitativamente, modesto, y desde su introducción se han
desarrollado otras ecuaciones empíricas mucho más exitosas (ecuación de Redlich--Kwong, ecuación de
Soave--Redlich--Kwong, etc.) para aplicaciones prácticas, que incluyen $5$ o más constantes
empíricas. Sin embargo, cualitativamente la ecuación de van der Waals es muy exitosa para predecir
las propiedades de los fluidos reales, incluyendo la llamada \emph{transición de fase
líquido--vapor} (es decir, el punto en el que el sistema pasa de gas a líquido).

Desde el punto de vista termodinámico no es posible obtener mayor información respecto a esta
ecuación, por lo que se debe recurrir a la mecánica estadística para obtener una derivación formal
de esta ecuación.

Para poder cerrar el sistema de ecuaciones de estado, se necesita una ecuación de estado térmica.
Una forma de hacerlo es a partir de experimentos; sin embargo, también se puede deducir esta
ecuación partiendo de la idea de obtener la ecuación térmica más simple y razonable que pueda ser
compatible con la ecuación de estado de van der Waals. No se puede simplemente tomar la ecuación de
estado del gas ideal, ya que el formalismo termodinámico impone una condición de consistencia entre
ambas, por lo que se puede esperar que, tal y como la ecuación mecánica tiene pequeñas diferencias
respecto al caso ideal, es de esperar que la ecuación térmica también tenga que ser ligeramente
modificada.

Si se escribe la ecuación de van der Waals como
\begin{equation*}
  \frac{P}{T}=\frac{R}{v-b}-\frac{a}{v^{2}}\,\frac{1}{T}
\end{equation*}
y buscamos una ecuación de estado térmica de la forma $\dfrac{1}{T}=f(u,v)$, ambas ecuaciones se
pueden integrar en la relación molar del cambio de entropía
\begin{equation*}
  \dd s=\frac{1}{T}\,\dd u+\frac{P}{T}\,\dd v
\end{equation*}
y obtener la ecuación fundamental $s=s(u,v)$. Sin embargo, esta relación requiere que, para que
$\dd s$ sea un diferencial perfecto (es decir, no requiere un factor integrante), se debe cumplir que
las derivadas cruzadas de segundo orden sean iguales, es decir,
\begin{equation*}
  \boxed{\frac{\partial^{2}s}{\partial u\,\partial v}=\frac{\partial^{2}s}{\partial v\,\partial u}}\qquad\text{(teorema de Schwarz)}.
\end{equation*}

El teorema de Schwarz establece las condiciones de exactitud para que un diferencial sea completo o
exacto.

\begin{teorema}[Schwarz]
  Sean $M,N$ funciones de clase $C^{1}$ en un dominio $D\subset\R^{2}$ simplemente conexo. Existe
  una función $f$ (de clase $C^{2}$) tal que
  \begin{equation*}
    \dd f=M(x,y)\,\dd x+N(x,y)\,\dd y,\qquad M=\pdv{f}{x},\quad N=\pdv{f}{y},
  \end{equation*}
  si y solo si
  \begin{equation*}
    \pdv{M}{y}=\pdv{N}{x},\qquad\text{es decir,}\qquad\frac{\partial^{2}f}{\partial x\,\partial y}=\frac{\partial^{2}f}{\partial y\,\partial x}.
  \end{equation*}
  Cuando esto ocurre, el diferencial es cerrado y, por lo tanto, exacto en un dominio simplemente
  conexo.
\end{teorema}
% NOTA: el manuscrito enuncia esto para f de clase C^2 con df=M dx+N dy y afirma "se cumple si y solo si dM/dy=dN/dx". Se reformuló con las hipótesis explícitas (M,N de clase C^1, dominio simplemente conexo) para que el enunciado sea correcto; es el mismo resultado.

Regresando a nuestro caso, se debe cumplir entonces que
\begin{equation*}
  \frac{\partial^{2}s}{\partial v\,\partial u}=\frac{\partial^{2}s}{\partial u\,\partial v}
  \qquad\text{o}\qquad
  \pdv{v}\left(\pdvc{s}{u}{v}\right)_{u}=\pdv{u}\left(\pdvc{s}{v}{u}\right)_{v},
\end{equation*}
es decir,
\begin{equation*}
  \left[\pdv{v}\left(\frac{1}{T}\right)\right]_{u}=\left[\pdv{u}\left(\frac{P}{T}\right)\right]_{v};
\end{equation*}
pero sabemos que $\dfrac{P}{T}=\dfrac{R}{v-b}-\dfrac{a}{v^{2}}\dfrac{1}{T}$; entonces, en este caso,
solo $\dfrac{1}{T}=\dfrac{1}{T}(u,v)$ depende de $u$, por lo que
\begin{equation*}
  \left[\pdv{v}\left(\frac{1}{T}\right)\right]_{u}=\left[\pdv{u}\left(\frac{R}{v-b}-\frac{a}{v^{2}}\frac{1}{T}\right)\right]_{v}
  =-\frac{a}{v^{2}}\left[\pdv{u}\left(\frac{1}{T}\right)\right]_{v},
\end{equation*}
y despejando $v^{2}$,
\begin{equation*}
  v^{2}\left[\pdv{v}\left(\frac{1}{T}\right)\right]_{u}=-a\left[\pdv{u}\left(\frac{1}{T}\right)\right]_{v}.
\end{equation*}
Si escribimos $\beta=\dfrac{1}{v}$ (es decir, $v=1/\beta$) y $\alpha=\dfrac{u}{a}$ (es decir,
$u=a\alpha$), aplicamos la regla de la cadena para cada derivada:
\begin{equation*}
  \pdv{v}=\pdv{\beta}{v}\pdv{\beta}=-\frac{1}{v^{2}}\pdv{\beta},
  \qquad
  \pdv{u}=\pdv{\alpha}{u}\pdv{\alpha}=\frac{1}{a}\pdv{\alpha}.
\end{equation*}
Por lo tanto
\begin{equation*}
  v^{2}\left[-\frac{1}{v^{2}}\pdv{\beta}\left(\frac{1}{T}\right)\right]_{u}=-a\cdot\frac{1}{a}\left[\pdv{\alpha}\left(\frac{1}{T}\right)\right]_{v}
  \quad\Longrightarrow\quad
  \left[\pdv{(1/v)}\left(\frac{1}{T}\right)\right]_{u}=\left[\pdv{(u/a)}\left(\frac{1}{T}\right)\right]_{v}.
\end{equation*}
Por lo que $\dfrac{1}{T}$ debe ser función de $1/v$ y $u/a$ de forma tal que ambas derivadas deben
ser iguales de acuerdo con el teorema de Schwarz; una posible forma es que $1/T$ dependa de la suma
$\left(\dfrac{1}{v}+\dfrac{u}{a}\right)$. Recordando que para el caso de gas ideal se tiene
$1/T=cR/u$, el cambio más simple posible es
\begin{equation*}
  \boxed{\frac{1}{T}=\frac{cR}{u+a/v}}
\end{equation*}

\begin{definicion}[Fluido ideal de van der Waals]
  El sistema caracterizado por las ecuaciones de estado para $P/T$ y para $1/T$,
  \begin{equation*}
    \frac{P}{T}=\frac{R}{v-b}-\frac{a}{v^{2}}\frac{1}{T},\qquad\frac{1}{T}=\frac{cR}{u+a/v},
  \end{equation*}
  es llamado \textbf{fluido ideal de van der Waals}.
\end{definicion}

Escribir $\dfrac{P}{T}=\dfrac{R}{v-b}-\dfrac{a}{v^{2}}\dfrac{1}{T}$ no es una forma adecuada de la
ecuación de estado, ya que $P/T$ depende de $1/T$; sustituyendo la ecuación de estado térmica se
puede encontrar la ecuación correcta para $P/T$:
\begin{equation*}
  \boxed{\frac{P}{T}=\frac{R}{v-b}-\frac{acR}{uv^{2}+av}}
\end{equation*}
Con estas ecuaciones es posible encontrar la relación fundamental en la representación entrópica,
\begin{equation*}
  \boxed{S=NR\ln\left[(v-b)\left(u+\frac{a}{v}\right)^{c}\right]+Ns_0}
\end{equation*}
donde $s_0$ es una constante. Esta relación, como el caso de la ecuación del gas ideal, no cumple
con el postulado de Nernst y no es válida para temperaturas bajas.

Algunos ejemplos de los valores de las constantes para el fluido ideal de van der Waals:
\begin{center}
  \begin{tabular}{lccc}
    \hline
    Gas & $a\ \left(\mathrm{Pa\,m^{6}/mol^{2}}\right)$ & $b\ \left(10^{-6}\,\mathrm{m^{3}/mol}\right)$ & $c$\\
    \hline
    He & $0.00346$ & $23.7$ & $3/2$\\
    Ne & $0.0215$ & $17.5$ & $3/2$\\
    $\mathrm{H_2}$ & $0.0248$ & $26.6$ & $5/2$\\
    Ar & $0.132$ & $30.2$ & $3/2$\\
    $\mathrm{N_2}$ & $0.136$ & $38.5$ & $5/2$\\
    CO & $0.151$ & $39.9$ & $5/2$\\
    $\mathrm{H_2O}$ & $0.544$ & $30.5$ & $3.1$\\
    $\mathrm{N_2O}$ & $0.384$ & $44.2$ & $7/2$\\
    \hline
  \end{tabular}
\end{center}
% NOTA: el manuscrito rotula las unidades de a como "Pa m^6" y de b como "10^{-6} m^3"; con v molar (m^3/mol), las unidades consistentes con P=RT/(v-b)-a/v^2 son Pa m^6/mol^2 y m^3/mol, que son las que se usan aquí. El gas "A" del manuscrito es el argón (Ar). Los valores 0.132 (Ar) y 3.1 (H2O) son los transcritos del manuscrito; verificar.

Estas constantes se obtienen al ajustar empíricamente las isotermas de van der Waals en la vecindad de
$273\ \mathrm{K}$; no son valores representativos para temperaturas más bajas.

\subsection{Ecuación de Clausius--Clapeyron}

% Pendiente: aún no está en las notas.

\subsection{Regla de fases de Gibbs}

% Pendiente: aún no está en las notas.

\section{Diagrama de fase de una sustancia con propiedades magnéticas. Ley de Curie}

% Pendiente: aún no está en las notas.

\section{Sólidos: ley de Dulong y Petit}

% Pendiente: aún no está en las notas.

\section{Reacciones químicas endotérmicas y exotérmicas}

% Pendiente: aún no está en las notas.

